\chapter{Running fBlockKit}
\label{ch:running}

\section{The menu}

Start the program (Section~\ref{ch:installation}) and you get a numbered menu.
Type the number of a function and press Enter. Every prompt tells you what it
wants and which value is the default; pressing Enter alone accepts the default.
At any menu, typing \code{0} leaves the program.

The menu numbers are part of the public interface: they are append-only and
never reordered, because a session script replays by number. In this manual the
chapter of a menu number documents that function; a cross-reference such as
\menu{9} points to it.

\begin{transcript}[title={The program starts with its menu (excerpt)}]
fBlockKit -- the f-block calculation toolkit (input generation + characterisation analysis; it runs no calculations)

Available functions:
   1  Check-up and characterisation (read an ORCA output -> analysis + diagnosis report)
   2  Coordination geometry and symmetry hints (read an XYZ structure)
   ...
  11  Point-charge crystal-field estimate (read an XYZ structure + charges)
   0  Quit
Choose a number (Enter = leave empty)>
\end{transcript}

\section{The interaction is a script}

Every thing you type at a prompt is one line of a script. A script is a text
file with one input per line:

\begin{itemize}
  \item an \textbf{empty line} means ``just press Enter'' (accept the default);
  \item a line starting with \code{\#} is a comment and is ignored;
  \item everything else is passed to the prompt exactly as if you had typed it.
\end{itemize}

You can obtain a script in two ways: record while you work
(\cmd{--record FILE}), or save the session so far with \menu{8}. Both record
the answer sequence itself; the act of saving is not part of the script, so
replaying it never asks for a save path.

A script replays with

\begin{codeblock}
fblockkit run script.txt
\end{codeblock}

The program prints exactly what it would have printed interactively, including
the prompts and your answers, so a replay is a complete transcript. Replaying
the same script twice produces \emph{byte-identical} output -- the program
writes no timestamps and nothing else unstable. This property is also the
project's regression acceptance criterion: the test suite replays the scripts
in this manual and compares them byte for byte.

\begin{notebox}[Reading a transcript]
The transcript boxes in this manual show both sides of the conversation: the
prompt line ends with \code{>} followed by the typed answer. An empty line in
a script appears as an empty answer.
\end{notebox}

\subsection{A worked recording}

\lstinputlisting[style=fbkcode,
  title={The script used for the recording example: \file{examples/scripts/m8\_script.txt}}]
  {examples/scripts/m8_script.txt}

The session performs two analyses and then saves its input at \menu{8}:

\lstinputlisting[style=fbkcode, firstline=60, lastline=67, title={Session excerpt (menu 8)}]
  {examples/expected/m8_script.txt}

The saved file contains exactly the six input lines typed \emph{before} the
saving action:

\lstinputlisting[style=fbkcode,
  title={The recorded script (\file{examples/expected/saved\_session.txt})}]
  {examples/expected/saved_session.txt}

Replaying that file reproduces the two analyses; when the replay reaches
\menu{8} the script has no further line, so the program reports the finished
input and leaves cleanly:

\lstinputlisting[style=fbkcode, firstline=63, lastline=67, title={Replay of the saved script}]
  {examples/expected/m8_replay.txt}

\section{The command line}

\begin{codeblock}
fblockkit                      interactive menu
fblockkit --record FILE        the same, recording your input as a script
fblockkit run SCRIPT.txt       replay a script
fblockkit search KEYWORDS      search the tool index (Chapter~\ref{ch:menu5})
fblockkit guide TOOL_ID        tool onboarding notes (Chapter~\ref{ch:menu6})
\end{codeblock}

Exit codes: \code{0} for a normal run; \code{2} for a usage problem (a missing
script file, an empty keyword, an unknown tool id). The program prints
diagnostics to standard output only -- there is no log file, and no state is
kept between runs.

\section{What gets written, and where}

\fbk{} writes new files \emph{next to the input file you gave it} (never
overwriting anything):

\begin{center}
\begin{tabularx}{\textwidth}{@{}lX@{}}
\toprule
Product & Written by \\
\midrule
\file{name.fbk.md}, \file{name.fbk.json} & \menu{1} (report), \menu{2}, \menu{10}, \menu{11} \\
\file{structure.fbk.inp} & \menu{3} (generated ORCA input) \\
\file{input.fix\_\textit{variant}.inp} & \menu{9} (corrected inputs) \\
\bottomrule
\end{tabularx}
\end{center}

Reports and generated files contain no timestamps, so a re-run overwrites the
previous product with identical content rather than creating clutter.
